\chapter{催眠の脱神話化：二つの作動原理とサイコ・サイバネティクス}
\epigraph{%
  \itshape ``Erickson was not an occultist wielding supernatural powers;
  he was a relentless strategist of human interaction.
  What onlookers mistook for mystical trance was merely the visible consequence
  of rigorous control over communication levels,
  leaving the subject with no behavioral exit except compliance.''
}{%
  \textbf{Jay Haley}\\
  \textit{Uncommon Therapy: The Psychiatric Techniques of Milton H. Erickson, M.D.}
}
催眠という現象について、世間の大半の人間、そして表層的なNLP学習者すらも致命的な誤解に囚われている。
彼らは催眠を「何か特殊な魔力や超越的な意識状態によって、相手を意のままに操る術」であると信じ込んでいる。
しかし、サイコ・サイバネティクスシステムの目標値（セットポイント）が\textbf{「確信 $\Omega^{\mathrm{ego}}_{\mathrm{predict:profit}}$」}であったことを思い起こしてほしい。
生体コンピュータは、外部から注入された目標値 $\Omega$ と現在の状態との誤差を最小化するように、全身の神経系と生理反応を自動制御（サーボ追従）する閉ループ系である。

この物理的本質を直視するならば、催眠と呼ばれる現象は神秘のベールを完全に脱ぎ捨て、以下の\textbf{二つの独立した工学的作動原理}へと還元される。

\section{カリスマ性による前提セット（サイコ・サイバネティクス追従）}

第一のパターンは、\textbf{「術者のカリスマ性というエビデンス群によって、被験者の内部に前提となる確信 $\Omega_{\mathrm{authority}}$ が強固にセットされ、その目標値へ向けてサイコ・サイバネティクスが合目的的に自己組織化する場合」}である。

被験者が「この高名なマスターは本物である」「この人物には私の心が見通されている」「彼なら私を劇的に変容させてくれるに違いない」という強い事前予期を抱いているとき、術者の堂々たる振る舞い、演出、場の熱気、過去の実績といった無数の外的シグナル群 $\langle V_e, A_e, K_e, O_e, G_e \rangle$ が、反論の余地のないエビデンスとして被験者の神経系疑似ベイズ推論機に雪崩れ込む。
その結果、次のマクロな確信が絶対的なセットポイントとして確定する。

\begin{equation}
    \boldsymbol{e}_{\mathrm{authority}} \nbuc{\hat{\Omega}}{\hat{\Omega}'} \Omega_{\mathrm{myth}}
\end{equation}
ここで $\Omega_{\mathrm{myth}} \equiv \text{「この術者の介入によって、私は変化する／従う」}$ という前提仮説の確定を意味する。

ひとたびこの「神話」自体が前提の確信 $\Omega_{\mathrm{myth}}$ として固定されると、サイコ・サイバネティクスのサーボ機構のスイッチが入る。
被験者の神経系は、\textbf{「その確信を現実に成就させるため」に、無意識下で全身の感覚受容、自律神経系、筋出力を合目的的に書き換えてしまう}のである。

\subsection{アイ・アクセシング・キューの虚構とバンドラーの悪戯}
この力学を象徴するのが、NLP教本に必ず登場する「アイ・アクセシング・キュー（視線方向による内的表象の特定理論）」である。
現代の認知科学および神経科学の追試において、視線の動きと内的表象（視覚・聴覚・体感覚）の間に普遍的・神経学的な因果律など存在しないことは、すでに周知の事実である。

では、なぜ実際のセッションにおいて、被験者はあたかも理論通りに視線を動かし、術者の指摘に驚愕するのか。
それこそが、前提としてセットされた確信 $\Omega_{\mathrm{myth}}$ の威力にほかならない。
被験者の脳は、「この術者は私の視線から思考を完全に読み取っている」という前提 $\Omega$ を維持するために、無意識下でその期待（セットポイント）に適合するよう自らの眼球運動を合目的的に同期させてしまうのである。

リチャード・バンドラーの極めて悪戯好き（playful）でトリックスター的な天才性は、まさにここにある。
彼は、神経科学的な根拠など最初から1ミリも存在しない「アイ・アクセシング・キュー」のような小道具を、\textbf{「私はあなたの心を完全に読んでいる」という前提の確信 $\Omega$ を相手の脳内へ強烈にインストールするための舞台装置（催眠の前提セット）として平然と使い倒していた}のだ。

NLPのテキストに散りばめられた一見もっともらしい「科学的風貌の理論」の多くは、客観的真理などではない。
被験者の批判的検閲（CEN）を納得させ、反論の余地なく $\Omega_{\mathrm{myth}}$ を確定させるために仕掛けられた、彼の卓越した催眠術師としての壮大な「ハッキング・コード（前提の埋め込み）」だったのである。

\section{不可知 $\UU$ の投入と解の保留（アルタードステート $\State(\mathrm{ALTER})$）}

催眠のもう一つの独立した作動原理は、術者のカリスマ性や事前のラポールすら本質的には必要としない。
それは、\textbf{「神経系疑似ベイズ推論機の収束を物理的に阻止し、解の確定を強制保留（$\boldsymbol{e}_{\mathrm{seq}} \nbuc{\hat{\Omega}}{\emptyset} \mid$）させることによって引き起こされるアルタードステート（変性意識状態）」}である。

\subsection{超限不可知 $\aleph$ による推論機のフリーズ}
生体の観測・計算能力を超え、その発生理由や因果構造を既存のフレームワークで了解できない現象――本稿で\textbf{「Unlimited Unknown（超限の不可知：$\aleph$）」}と定義した入力シグナル――が注入されたとき、神経系は解の出ないTOTEループへと突入する。

\begin{equation}
    \aleph \nbuc{\hat{\Omega}}{\emptyset} \mid \quad \implies \quad \State(\mathrm{ALTER})
\end{equation}

ダブルバインド、急激な文脈の破断（パターン・インタラプト）、ミルトン・モデルによる極限の抽象化などにより、エゴの利害（profit / threat）を激しく刺激しながらも、いかなるナラティブ解 $\Omega$ にも収束できない矛盾シグナル $\aleph$ を連続投入する。
エゴにとって切迫した状況であるため計算を打ち切る（Exitする）ことが許されず、しかし提示される情報が既存のどの $\Omega^{\mathrm{ego}}$ にも適合しないとき、通常の思考ルーチンは過負荷によって完全にクラッシュする。

\subsection{全検索ゲートの開放と被暗示性の極大化}
解が確定しない空転状態において、脳は計算を放棄するのではなく、むしろ逆の挙動を示す。
確信 $\Omega$ という「確定した枠組み（解釈のゲシュタルト）」が成立しないため、神経系疑似ベイズ推論機は解を見出そうと、\textbf{エピソード記憶、身体感覚、無意識の連合ネットワークを含むあらゆる神経系の全領域を総動員し、広範かつ並列的な仮説探索を激しく駆動させ続ける}。

このとき生じるのが、変性意識状態特有の神経力学である。
\begin{itemize}
    \item \textbf{DMNによる固定ナラティブの停止}：\\
    自己防衛の既成ストーリー（$Ad^{\mathrm{ego}}$）がフリーズし、自己と他者、内面と外界の境界認識が一時的に融解する。
    \item \textbf{被暗示性（Suggestibility）の極大化}：\\
    脳があらゆる回路を開いて猛烈に「解（確信）」を欲している飢餓状態にあるため、外部から投入される微小な言語シグナルや象徴的介入が、エゴの検閲（CEN）を受けることなく、新たな確信 $\Omega^{\mathrm{wisdom}}$ として直接インストールされる。
\end{itemize}

\subsection{ソムナンビュリストの計算論的解体：トラウマ的予期 $\hat{\Omega}_{\mathrm{trauma}}$ と $\mathcal{S}(\mathrm{ALTER})$ への相転移}
%FIXME: 記号が間違っている。整合性が着くように直すこと

現代精神医学の臨床現場において、「ソムナンビュリスト（Somnambulist）」という語は、単なる睡眠時随伴症の一種である「夢遊病者」を指す狭隘な病名へと退廃してしまった。これは、患者個々の認知アーキテクチャの因果連鎖を追尾することを放棄し、統計的再現性とチェックリスト方式による症候分類（DSM等）へと矮小化させた現代精神医学の構造的怠惰であり、さらには対症療法的な投薬プロトコルによって利益を固定化する薬物利権の帰結にほかならない。

しかし、ピエール・ジャネ（Pierre Janet）やアンブロワーズ＝オーギュスト・リエボー（Ambroise-Auguste Liébeault）らが探求した古典催眠の文脈において、ソムナンビュリストとは「極めて高い催眠感受性を持ち、瞬時に深深度のトランスへと沈降する特異な被験者」を指す術語であった。この現象の物理的・情報論的実態を、生体疑似ベイズ推論機およびサイコ・サイバネティクスのフレームワークによって解剖するならば、それは神秘的な体質などではなく、\textbf{「強烈かつ未解決なトラウマ的事前予期 $\hat{\Omega}_{\mathrm{trauma}}$ を内部に抱え込んでいる状態」}と等価であると再定義される。

\begin{figure}[htbp]
\centering
\begin{tikzpicture}[node distance=1.8cm, auto]
    \node [draw, rectangle, rounded corners] (input) {環境ノイズ $e_{\mathrm{noise}}$};
    \node [draw, rectangle, fill=gray!10, right of=input, xshift=2.5cm] (threat) {$\Omega^{\mathrm{ego}}_{\mathrm{predict:threat(trauma)}}$ (確信の凍結)};
    \node [draw, diamond, aspect=2, below of=threat, yshift=-0.5cm] (cause) {$\Omega^{\mathrm{ego}}_{\mathrm{cause}}$ 確定?};
    \node [draw, rectangle, fill=red!10, left of=cause, xshift=-2.5cm] (alter) {$\mathcal{S}(\mathrm{ALTER})$ (全検索ゲート開放・超警戒)};
    
    \draw [->, thick] (input) -- node {$\hat{\Omega}_{\mathrm{trauma}}$} (threat);
    \draw [->, thick] (threat) -- (cause);
    \draw [->, thick] (cause) -- node [above] {No ($\emptyset |$)} (alter);
    \draw [->, dashed] (alter) |- (input);
\end{tikzpicture}
\caption{ソムナンビュリストにおけるトラウマ的予期と $\mathcal{S}(\mathrm{ALTER})$ への遷移ループ}
\label{fig:somnambulist_loop}
\end{figure}

ソムナンビュリストの神経系内部では、過去の強烈な情動負荷（$\Lambda_{\mathrm{terror}}$）によって刻印されたトラウマ的予期 $\hat{\Omega}_{\mathrm{trauma}}$ が、環境からのわずかなノイズ刺激 $e_{\mathrm{noise}}$ を契機として、極めて高い剛性を持った脅威の未来予測確信へと一瞬で相転移する。
\begin{equation}
e_{\mathrm{noise}} : \frac{\hat{\Omega}_{\mathrm{trauma}}}{\hat{\Omega}'_{\mathrm{trauma}}} \Longrightarrow \Omega^{\mathrm{ego}}_{\mathrm{predict:threat(trauma)}}
\end{equation}

ここで決定的な力学的分岐が生じる。サイコ・サイバネティクス制御器がこの壊滅的脅威 $\Omega^{\mathrm{ego}}_{\mathrm{predict:threat(trauma)}}$ を回避・遮断するためには、「なぜその脅威が発生したのか」という因果帰属の確信 $\Omega^{\mathrm{ego}}_{\mathrm{cause}}$ が特定されていなければならない。しかし、トラウマの本質とは「理不尽かつ不可解な外力の直撃」であり、その原因同定は原理的に未解決（$\emptyset |$）のまま宙吊りにされている。

すなわち、生体は「確実に破滅的脅威が到来する」という絶対的確信に拘束されながら、同時に「その脅威がどこから、どのような因果で生じるのか（$\Omega^{\mathrm{ego}}_{\mathrm{cause}}$）」を計算できないという致命的な閉塞（デッドロック）に直面する。このとき、不確実性を極限まで嫌悪する DMN モジュールは計算を停止することを許されず、外部環境から脅威の兆候と解釈コードを貪欲かつ猛烈にサンプリングしようと試みる。

\begin{equation}
\left\{
\begin{aligned}
\Omega^{\mathrm{ego}}_{\mathrm{predict:threat(trauma)}} & \longrightarrow \text{確定 (Set-point 固定)} \\
\Omega^{\mathrm{ego}}_{\mathrm{cause}} & \longrightarrow \emptyset | \quad (\text{未解決・探索継続})
\end{aligned}
\right\}
\Longrightarrow \mathcal{S}(\mathrm{ALTER})
\end{equation}


この結果として引き起こされるのが、極度の緊張と注意の集中を伴いながら、全検索ゲートを全方位に開放した変性意識状態 $\mathcal{S}(\mathrm{ALTER})$ への強制突入である。

古典催眠においてソムナンビュリストが「術者の微小な暗示に瞬時に反応し、幻覚すら容易に現前化させる」と記録されてきた真の理由はここにある。彼らは従順なのでも、心が柔軟なのでもない。内奥に抱えた $\Omega^{\mathrm{ego}}_{\mathrm{predict:threat(trauma)}}$ の恐怖から逃れるため、飢餓状態に陥った推論機が、外部から投入される術者の言語シグナル（暗示）を「この脅威を説明し、ループを Exit させてくれる唯一の救済解」として、批判的検閲（CEN）を一切経由せずに渇望・過剰マッチングさせているのである。

ソムナンビュリズムとは、精神医学がラベルを貼って向精神薬で鈍麻させるべき「機能不全の病理」などではない。それは、未解決のトラウマ的破滅予測を抱えた生体コンピュータが、自己防衛のために外部情報へ全帯域を開放した、極限の適応的サバイバル・ステートなのである。

\subsection{「神話への服従」か「計算機の停止」か}
このように解剖したとき、世に言う催眠とは、
\begin{enumerate}[label=(\arabic*)]
    \item \textbf{「術者のカリスマ性をエビデンスとして $\Omega$ を強固に確定させ、サイコ・サイバネティクスにその神話を自動達成させるトップダウンの閉ループ追従」}であるか、あるいは、
    \item \textbf{「既存の仮説空間をショートさせ、$\Omega$ を確定不能（$\mid$）に陥れることで全神経系を探索モードへと強制開放する状態（アルタードステート）」}
\end{enumerate}
のいずれか――あるいはその巧妙な組み合わせにすぎない。

魔術師が起こしているのは奇跡ではない。
相手の脳という情報処理装置に対し、「絶対的なセットポイント $\Omega$ を与えて下位システムを自動駆動させる」か、「セットポイントの計算そのものをクラッシュさせて無防備な空白（$\mid$）を作り出す」かという、冷徹なシステム制御を実行しているにすぎないのである。


